۱۸۹  
«متوجه‌ای چه می‌گم،» میستری با شدت گفت. «می‌رم یه چاقو بگیرم و اون‌وقت می‌فهمیم دقیقاً کی باید توی راهرو باشه وقتی برگردم.»

میستری به‌سمت خانه قدم برداشت و بقیهٔ ما نگاهی نگران با هم رد و بدل کردیم. دوباره، رفتار او را از سفر جاده‌ای‌مان به یاد آوردم: به یاد مرزی افتادم که وقتی به او گفتم چه کار کند، ناگهان خشمگین شد و مسائل پدری‌اش بیدار شدند.

کشیش در را به شدت بست و من و کارولین در آشفتگی صبح پس از پست‌های پر از کدئین او، به آرامی فرار کردیم. میستری در صندلی عقب خودروی کارولین فرو افتاده بود، درون پتو پیچیده و کلاهش را به پایین کشیده بود تا چشمانش را بپوشاند. به‌جز درخواست که او را در کنادوی خانوادگی‌اش پیاده کنیم، حرفی نزد، که برای او نادر بود. مرا به یاد سفر جاده‌ای‌مان در اروپای شرقی می‌انداخت. اما این بار، میستری بیمار نبود—حداقل نه از نظر جسمی.

ماشین را پارک کردیم و با آسانسور به آپارتمان خواهرش در طبقه بیستم رفتیم. آپارتمانی دو خوابه و شلوغ با مردمی زیاد. مادر میستری، زنی آلمان‌تبار و سنگین، بر روی صندلی کهنه گل‌دار نشسته بود. خواهرش، مارتینا، با دو فرزندش و همسرش، گری، کنار او بر روی مبل جا گرفته بودند. پدر میستری در آپارتمانی چهار طبقه بالاتر به دلیل بیماری کبد ناشی از مصرف مادام‌العمر الکل بستری بود.

«پس چرا دوست‌دختری باهات نیست؟» شیلن، برادرزاده سیزده سالهٔ میستری، او را سرزنش کرد. او همه چیز دربارهٔ دختران میستری را می‌دانست. میستری اغلب از برادرزاده‌هایش برای نشان‌‌دادن جنبهٔ آسیب‌پذیر و پدرانه‌اش به زنان استفاده می‌کرد. او به برادرزاده‌هایش واقعاً عشق می‌ورزید و به نظر می‌رسید وقتی آنها را می‌دید، دوباره جان می‌گرفت.

شوهرخواهر میستری، گری، برایمان چند ترانه پاپی که ساخته بود، اجرا کرد. بهترین آنها ترانه‌ای به نام «فرزند کازانووا» بود که میستری با صدای نزدیک به کرکننده‌ای همراهی کرد. به نظر می‌رسید با شخصیت اصلی ترانه همذات‌پنداری می‌کند.

کارولین و من بعد از آنجا را ترک کردیم. دختران تا دم در آسانسور ما را دنبال کردند، خندان و فریادزنان، و با میستری به دنبال ما آمدند. ناگهان در باز شد و مردی با یقهٔ روحانی به دختران نگاهی محکم و معترضانه انداخت.

«شما نباید توی راهرو این همه سر و صدا کنید،» او گفت.

صورت میستری سرخ شد. «می‌خواهی چه کار کنی؟» او پرسید. «چون به نظر من باید بکنیم. این‌ها دختران جوانی هستند. دارند خوش می‌گذرانند.»

کشیش گفت، «خوب، آنها می‌توانند در جایی که ساکنان دیگر را مزاحم نمی‌کنند، خوش بگذرانند.»